% Storyline 3 (final), paragraph 2: the punishment dilemma in plain terms; knowledge versus revision.
The trouble lies in punishment. Stable norms allow people to refuse help to those in bad standing without losing their own good standing, because the refusal is seen as justified\cite{panchanathan2003,ohtsuki2004,okada2020dilemmas}. Whether a refusal looks justified, however, depends on the observer's own view of the person refused. An observer who still regards that person as good condemns the punisher, and disagreement about one individual spreads to those who act towards it\cite{hilbe2018,uchida2013,radzvilavicius2019}. Those who punish defectors thus put their own reputation at risk, a punishment dilemma that undermines the sanctions on which cooperation depends\cite{okada2020dilemmas,michelmata2024,pires2025}. The harm has two components, how often justified punishers are condemned and how much each condemnation changes their reputation. The first depends mainly on what observers know, the second on how they revise what they believe.
